% ===========================================================================
%  问题三正文（由 main.tex 以 \input 引入）
%  数字来源：
%    03_代码/p3_analysis.py    主结果、八种组合、口径敏感性、参数扫描、九项校验、预报精度
%    03_代码/p3_calibration.py 1 月标定面与样本外检验
%    03_代码/p3_tables.py      正文表 p3_tables.tex / 附录表 p3_tables_appendix.tex
%    03_代码/p3_figures.py     插图 04_图/p3_fig*.png
%  本文件不含任何手工转录的数字。表号由 LaTeX 自动编号，正文一律用 \ref。
% ===========================================================================

\section{问题三：多版本预报驱动的滚动购电与调整模型}

\subsection{模型框架：三个购电量与交付块策略}

问题三的核心困难不在物理模型——储能、能量平衡与效率约束与问题二完全一致——而在
\textbf{"计划"与"调整"这两个概念必须被严格区分}。为此首先定义三个互不相同的量：

\begin{itemize}
  \item $g_t^0$：当天 $0{:}00$ 依据当时唯一可得的 $0{:}00$ 预报确定的
        \textbf{初始计划购电量}，一旦确定即冻结，后续任何调整都不得覆盖它；
  \item $x_t$：该时段\textbf{最终生效的完整购电量}（不含紧急购电）。它是一个
        完整量，\textbf{不是}相对于 $g_t^0$ 的增量，\textbf{也不是}该时段被
        修订过的多个版本之和；
  \item $r_t$：实际运行中额外发生的\textbf{紧急购电量}。
\end{itemize}

于是全天计划量为 $G^{\mathrm{plan}}=\sum_tg_t^0$，最终生效常规量为
$G^{\mathrm{eff}}=\sum_tx_t$，紧急量为 $G^{\mathrm{emg}}=\sum_tr_t$，微网实际
\textbf{收到}的电量为 $G^{\mathrm{recv}}=G^{\mathrm{eff}}+G^{\mathrm{emg}}$——
\textbf{不能再把 $g_t^0$ 加一遍}。

由于 $6{:}00$、$12{:}00$、$18{:}00$ 都可以修改购电，同一个时段在理论上可能被改
多次，因而必须先约定结算方式。本文采用\textbf{交付块（delivery block）}策略：
把一天切成四个六小时交付块，每个块\textbf{只在它的块首被正式提交一次}：

\begin{center}
\footnotesize
\begin{tabular}{lcc}
  \toprule
  决策时刻 & 本时刻正式提交的交付块 & 时段索引 \\
  \midrule
  $0{:}00$  & $0{:}00$--$6{:}00$   & $0$--$35$ \\
  $6{:}00$  & $6{:}00$--$12{:}00$  & $36$--$71$ \\
  $12{:}00$ & $12{:}00$--$18{:}00$ & $72$--$107$ \\
  $18{:}00$ & $18{:}00$--$24{:}00$ & $108$--$143$ \\
  \bottomrule
\end{tabular}
\end{center}

每个十分钟时段因此只与自己的 $g_t^0$ 结算一次，"最终净偏差结算一次"与"逐次变更
累加收费"之间的歧义被彻底消除。$0{:}00$ 的求解对全部 $144$ 段给出计划，但只有第一
块不可再改；后续各次求解对未来时段仍会一并算出候选值，它们只是\textbf{展望值}：
不提交、不计费，也不写入结果文件中的调整购电量，到对应交付块开始时再用当时的新
信息重新决定。这是为消除结算歧义而选定的策略范围，而非题目强制每次只能调整六小时。

\subsection{预报数据的时间口径}
\label{sec:p3-align}

\textbf{整点功率还是小时平均？}附件 3 的每条记录给出"预报 1 小时"至"预报 24
小时"共 $24$ 个\textbf{功率}值。它究竟是目标整点的瞬时功率，还是目标小时内的平均
功率，决定了整点之间该如何取值。本文用对齐扫描判定：把预报值与附件 2 的实际光伏
按不同时间平移对齐，比较全天平均绝对偏差，在正式区间内得到整点对齐时
$191.51$~kW、前移一小时 $808.67$~kW、按小时平均 $268.02$~kW（正式区间
2025-02-01 至 2025-12-31 的十分钟平均绝对偏差）。扫描脚本与输出见支撑材料
\texttt{p3\_alignment\_scan.py} 与 \texttt{p3\_alignment\_scan.json}。整点口径的
偏差最小且差距显著，故\textbf{按整点功率处理}。

\textbf{首端锚点：插值的左端必须是已经实现的量。}发布时刻为
$\tau\in\{0,6,12,18\}$ 的版本给出目标 $\tau+1,\dots,\tau+24$ 共 $24$ 个整点功率。
要把它们铺到十分钟网格上，分段线性插值还缺一个左端节点；若用 $\tau$ 时刻自身的
预报充当左端，就等于假定了 $\tau$ 时刻的功率已知，而附件 3 恰恰\textbf{没有}提供
"预报 $0$ 小时"这一项。本文取左端锚点为发布时刻所在时段的\textbf{实际}光伏功率
（$0{:}00$ 版取 $0$，因为附件 2 中全年每天 $0{:}00$--$0{:}10$ 的实际光伏恒为零），
节点为 $(\tau,\tau+1,\dots,\tau+24)$，在相邻节点之间作分段线性插值，只在
$\tau<\text{时段中点}\le\tau+24$ 内取值，并截断到非负；区间之外的时段该版本本就
不预测，返回空值而\textbf{不填零}。锚点取自附件 2 中\textbf{已经发生}的实测功率，
整条曲线只依赖 $\tau$ 时刻已经掌握的信息。

\textbf{负载与发布时间管理。}负载没有外部预报，只用历史实际值因果外推：优先取
\textbf{上一周同一星期几}（$d-7$）的实测曲线，退化时取 $d-1$，再退化时取已有历史
均值。光伏则严格按"发布时间不晚于当前决策时刻"筛选：第 $\tau$ 次求解只读取
$\tau$ 时刻\textbf{已经发布}的那一版预报，因此 $0{:}00$ 绝不会看到当天 $6{:}00$
才发布的信息。

图~\ref{fig:p3-versions} 把四个发布版本对同一个日期的光伏预报与事后实际值画在一起：
$0{:}00$ 版覆盖全天，$6{:}00$、$12{:}00$、$18{:}00$ 版的覆盖时段依次缩短。覆盖越短
的版本对它所覆盖的那一段也越准，但\textbf{精度高不等于价值大}——这四个版本各自的
边际价值是本节第（2）条要回答的问题。

\begin{figure}[htbp]
  \centering
  \includegraphics[width=.66\textwidth]{p3_fig1_versions.png}
  \caption{同一目标日期的多版本光伏预报与事后实际值}
  \label{fig:p3-versions}
\end{figure}

\subsection{计费口径与调整模型}
\label{sec:p3-billing}

题目规定："计划购电量高于调整购电量的部分，违约电价是交易时刻电价的 $50\%$；
调整购电量高于计划购电量的部分，超出部分的电价是交易时刻电价的 $1.5$ 倍"，而
"总的购电费用包括计划购电费用、紧急购电费用和调整购电量的相关费用"。这里存在
一处必须\textbf{显式声明}的歧义：调减时被取消部分的\textbf{原价是否退回}，题面
并未说明。本文以两种口径分别求解，并在论文中\textbf{同时报告}。

\textbf{主口径（退款：取消部分退回原价，只收 $50\%$ 违约费）。}
$$
J^{\mathrm{refund}}
=\sum_t\Big[p_tg_t^0+1.5p_t(x_t-g_t^0)_+-0.5p_t(g_t^0-x_t)_++5p_tr_t\Big],
\qquad (y)_+=\max(y,0).
$$
其中 $1.5p_t$ 是上调部分的\textbf{总价}，不是"先付 $p_t$ 再加 $1.5p_t$"成为
$2.5p_t$；$-0.5p_t$ 项表示调减部分先退回原价 $p_t$、再按 $50\%$ 收违约费，净得
$0.5p_t$ 的贷方。该口径下\textbf{调减可以降低成本}，优化允许双向调整。

\textbf{敏感性口径（不退款：调减仍付原价，另收 $50\%$ 违约费）。}
$$
J^{\mathrm{norefund}}
=\sum_t\Big[p_tg_t^0+1.5p_t(x_t-g_t^0)_++0.5p_t(g_t^0-x_t)_++5p_tr_t\Big].
$$
此时调减既付原价又付违约费，"不调"总不比"调减"贵，最优解不会真正调减；用同一套
模型、同一套标定与执行规则重跑，得到的是\textbf{同一决策机制在该读法下的账本}，
而不是另一条策略。两种口径的分段线性化方式相同：引入非负的增量与减量变量
$u_t,v_t\ge0$，令
$$
x_t-g_t^0=u_t-v_t ,\qquad 0\le v_t\le g_t^0 ,
$$
则调价项为 $1.5p_tu_t-0.5p_tv_t$（主口径）或 $1.5p_tu_t+0.5p_tv_t$（不退款
口径），\textbf{都无需引入二元变量}——主口径下把 $u_t,v_t$ 同时抬高 $\delta$ 只会
平白多付 $p_t\delta$，最优解中二者不会同时为正。两臂的全年结果并列于
表~\ref{tab:p3-billing}，相差 $1.49\%$；\textbf{本文不对哪一种读法作对错判断}，
只报告差异的量级，并把它与预报本身的收益（第（2）条，$10.41\%$）并列衡量。

\textbf{为何这是"调整"而不是"重新购电"。}在调整时点，原计划费用已经确定，在目标
函数中属于常数，因此不再重复计入；已经执行过的时段产生的调整费与紧急费同样不再
优化，只保留在实际账本中。

\subsection{计划与调整两个线性规划}

\textbf{$0{:}00$ 初始计划。}沿用问题二的风险净负荷思想，以
$\widetilde N_t$ 为需求求解
$$
\min\ \sum_{t=0}^{143}p_tg_t^0+\lambda\xi ,
\qquad \xi\ge E^{\mathrm{tar}}-\bar E_{143},\ \ \xi\ge0 ,
$$
$$
g_t^0+\bar d_t=\widetilde N_t+\bar c_t+\bar w_t ,
\qquad
\bar E_t=\bar E_{t-1}+\eta_c\bar c_t-\frac{\bar d_t}{\eta_d} ,
\qquad
1200\le\bar E_t\le10800 ,
$$
$$
0\le\bar c_t\le M,\quad 0\le\bar d_t\le M ,
\qquad M=5000\times\tfrac16=\tfrac{2500}{3}\ \mathrm{kWh}.
$$
其中 $\bar c_t,\bar d_t,\bar E_t$ 是\textbf{参考储能轨迹}，$\bar w_t$ 是规划弃电，
\textbf{都不是实际结果}。该模型\textbf{不含紧急购电变量}：$0{:}00$ 只有该时点发布
的预报版本可用，没有更新的预报，也没有当日的实际观测，模型只按风险净负荷安排一个
"量入为出"的计划，缺口与盈余都交给后续的调整层与执行层。购电量 $g_t^0$ 只有非负
约束，其有效上限由平衡式与储能容量自然给出——多买的电只能进储能或弃掉，两者都要
付出代价，故最优解不会无界买电。

\textbf{发布时点的调整模型。}设 $\mathcal T_\tau$ 为当前时刻之后尚未执行的时段，
以 $x_t$ 为候选完整购电量：
$$
\min\ \sum_{t\in\mathcal T_\tau}\big(1.5p_tu_t-0.5p_tv_t\big)
+\sum_{t\in\mathcal T_\tau}5p_t\bar r_t+\lambda\xi ,
$$
$$
x_t=g_t^0+u_t-v_t ,\qquad 0\le v_t\le g_t^0 ,
\qquad
x_t+\bar r_t+\bar d_t=\widetilde N_t^{(\tau)}+\bar c_t+\bar w_t ,
$$
储能递推式与上下界、充放电上限与计划模型完全相同。

两个模型都是\textbf{线性规划}而非混合整数规划：充放电互斥不需要二元变量来强制，
因为在 $\eta_c\eta_d=0.81<1$ 的效率结构下，同一时段同时充放电会净损失能量，任何
最优解都不会这么做。这里有两个容易被忽略的实现要点。其一，每次求解必须以
\textbf{当前实际储电量}为初值，而不是上一轮的预测电量——这是"滚动"的全部意义所在。
其二，$\widetilde N_t^{(\tau)}$ 与 $g_t^0$ 使用\textbf{不同的}信息集：前者用最新
发布的预报和最新误差分位数算出，后者在 $0{:}00$ 就冻结。更新预报后，即便不改变
$x_t$ 也可以重排储能、使未来缺口下降，因此代码在每次求解后\textbf{整体刷新参考
轨迹 $\bar E_t$}，但\textbf{只把当前交付块}的 $x_t$ 写回生效购电量，其余全部丢弃。

\textbf{执行层。}每十分钟的实际充放电与紧急补购完全复用问题二的因果规则——这种
"每个决策时刻只依据当前已实现状态重排未来轨迹"的做法正是模型预测控制的基本思想
\cite{ref:mpc}。令使用储能前的盈余为
$b_t=x_t+v_t^{\mathrm{act}}-\ell_t^{\mathrm{act}}$，用当前参考轨迹构造储备线
$R_t=\mathrm{clip}\big(1200+\rho(\bar E_t-1200),\,1200,\,10800\big)$，则
$$
c_t=\min\Big\{(b_t)_+,\ M,\ \frac{10800-E_{t-1}}{\eta_c}\Big\},
\qquad
d_t=\min\Big\{(-b_t)_+,\ M,\ \eta_d(E_{t-1}-R_t)_+\Big\},
$$
$$
r_t=(-b_t-d_t)_+ ,\qquad w_t=(b_t-c_t)_+ ,\qquad
E_t=E_{t-1}+\eta_cc_t-\frac{d_t}{\eta_d},
$$
并满足逐段守恒 $x_t+r_t+v_t^{\mathrm{act}}+d_t=\ell_t^{\mathrm{act}}+c_t+w_t$。
规则\textbf{只依赖当前与历史信息}：$b_t$ 用到当段实际负载与光伏，$E_{t-1}$ 是
上一段的实际状态。当日实际日末电量直接作为次日初值，不要求日循环；$1$ 月作为预热
期，$2$ 月 $1$ 日的初值由 $1$ 月的实际运行得到。

\subsection{参数标定：1 月联合标定与样本外检验}

需要标定的参数有四个：风险分位数 $\alpha$、储备线比例 $\rho$、日末储备引导权重
$\lambda$ 与分位数窗口 $W$。题目要求"每天 $0{:}00$ 制定当天计划购电策略"，因此
任何参数都只能由\textbf{已经过去}的数据决定。本文把 2025 年 1 月当作预热月，在这
$31$ 天上做一次联合标定，把选出的参数\textbf{冻结}，此后 $2$ 月 $1$ 日--$12$ 月
$31$ 日的评价区间内\textbf{不再回头修改}。

标定面为 $\alpha\in\{0.4,\dots,0.8\}\times\rho\in\{0.35,0.5,0.7\}
\times W\in\{14,21,28\}$ 的 $45$ 点网格，每组参数在 1 月上完整回放一次，取
\textbf{含 $5$ 倍紧急费用的实际总费用}最小者；$\lambda$ 另行单独扫描。结果见
表~\ref{tab:p3-calib}：标定面在 $\alpha$ 方向有明确的极小点 $\alpha=0.6$，在
$\rho$ 与 $W$ 方向则相当平坦。

\textbf{为什么取标定面的中心取值，而不是它的最小点。}1 月标定面上的最优取值是
$(\alpha,\rho,W)=(0.6,0.35,14)$，1 月费用 $1{,}763{,}097$~元，比本文采用的
$(0.6,0.5,28)$ 的 $1{,}772{,}822$~元省 $9{,}725$~元。但把这两组参数分别拿去跑
完整评价区间，前者 $13{,}461{,}840$~元、后者 $13{,}418{,}626$~元——\textbf{在 1
月上更优的那一组反而多花 $43{,}213$~元}。这说明 $31$ 天的标定窗口已经足以把标定
面上的细小凸起拟合进去，直接取网格最小点就是过拟合；因此本文取标定面的
\textbf{中心取值} $(\alpha,\rho,\lambda,W)=(0.6,0.5,1.0,28)$，并在
表~\ref{tab:p3-param} 中报告单参数扰动的完整灵敏度。$\lambda$ 取 $1.0$ 而非 1 月
最优的 $0.5$：两者在评价区间上只差 $9{,}852$~元（$0.07\%$），但 $\lambda=0$ 会把
日末储电量抽到 $1{,}806$~kWh（下界为 $1200$），属于把成本推给次日的透支行为，故
保留非零引导权重。

题目未给出获取预报本身的费用，因此\textbf{不应}自行增加"预报购买费"；"是否需要
其他时刻的预报"只能表达为：某次更新是否降低实际总费用、改善是否跨月份稳定、以及
能否被其他更新替代。

% 本文件由 03_代码/p3_tables.py 自动生成，请勿手工修改。
% 数据源：06_支撑材料/p3_results.json

\begin{table}[htbp]
  \centering
  \caption{指定日期六个时段的初始计划购电量 $g^0_t$、最终生效购电量 $x_t$、正式调整量与紧急购电量（单位：kWh）}
  \label{tab:p3-buy}
  \scriptsize
  \begin{tabular}{llrrrr}
    \toprule
    日期 & 时段 & 初始计划 $g^0_t$ & 最终生效 $x_t$ & 调整量 & 紧急购电 $r_t$ \\
    \midrule
    \multirow{6}{*}{2025.3.20} & 10{:}00-10{:}10 & 0.00 & 0.00 & 0.00 & 0.00 \\
     & 12{:}00-12{:}10 & 509.51 & 509.51 & 0.00 & 0.00 \\
     & 14{:}00-14{:}10 & 0.00 & 0.00 & 0.00 & 0.00 \\
     & 16{:}00-16{:}10 & 647.39 & 647.39 & 0.00 & 0.00 \\
     & 18{:}00-18{:}10 & 698.72 & 704.31 & 5.59 & 0.00 \\
     & 20{:}00-20{:}10 & 0.00 & 0.00 & 0.00 & 0.00 \\
    \midrule
    \multirow{6}{*}{2025.6.21} & 10{:}00-10{:}10 & 0.00 & 0.00 & 0.00 & 0.00 \\
     & 12{:}00-12{:}10 & 0.00 & 0.00 & 0.00 & 0.00 \\
     & 14{:}00-14{:}10 & 0.00 & 0.00 & 0.00 & 0.00 \\
     & 16{:}00-16{:}10 & 241.46 & 241.46 & 0.00 & 0.00 \\
     & 18{:}00-18{:}10 & 0.00 & 0.00 & 0.00 & 0.00 \\
     & 20{:}00-20{:}10 & 0.00 & 0.00 & 0.00 & 0.00 \\
    \midrule
    \multirow{6}{*}{2025.9.23} & 10{:}00-10{:}10 & 0.00 & 0.00 & 0.00 & 0.00 \\
     & 12{:}00-12{:}10 & 309.90 & 309.90 & 0.00 & 0.00 \\
     & 14{:}00-14{:}10 & 0.00 & 0.00 & 0.00 & 0.00 \\
     & 16{:}00-16{:}10 & 526.15 & 526.15 & 0.00 & 0.00 \\
     & 18{:}00-18{:}10 & 755.68 & 781.52 & 25.84 & 0.00 \\
     & 20{:}00-20{:}10 & 0.00 & 0.00 & 0.00 & 0.00 \\
    \midrule
    \multirow{6}{*}{2025.12.21} & 10{:}00-10{:}10 & 0.00 & 0.00 & 0.00 & 0.00 \\
     & 12{:}00-12{:}10 & 979.36 & 979.36 & 0.00 & 0.00 \\
     & 14{:}00-14{:}10 & 0.00 & 0.00 & 0.00 & 0.00 \\
     & 16{:}00-16{:}10 & 852.80 & 852.80 & 0.00 & 0.00 \\
     & 18{:}00-18{:}10 & 733.72 & 733.72 & 0.00 & 0.00 \\
     & 20{:}00-20{:}10 & 0.00 & 0.00 & 0.00 & 0.00 \\
    \bottomrule
  \end{tabular}
  \par\vspace{2pt}
  \begin{minipage}{.92\textwidth}\footnotesize
    续表：全天费用分解（单位：元）。计划费按 $g^0_t$ 全额收取，调整费按 $1.5p_t(x_t-g^0_t)$ 另计，二者与紧急费之和即当日实际总费用。
  \end{minipage}
  \par\vspace{4pt}
  \begin{tabular}{lrrrrrr}
    \toprule
    日期 & 初始计划量 & 最终常规量 & 紧急购电量 & 计划费 & 调整费 & 当日合计费用 \\
    \midrule
    2025.3.20 & 65,739.8 & 67,183.3 & 371.0 & 40,069 & 1,744 & \textbf{43,632} \\
    2025.6.21 & 35,596.8 & 35,608.9 & 0.0 & 20,694 & 20 & \textbf{20,714} \\
    2025.9.23 & 65,479.3 & 66,526.2 & 673.8 & 40,323 & 1,512 & \textbf{45,683} \\
    2025.12.21 & 96,942.5 & 96,942.5 & 12.1 & 62,044 & 0 & \textbf{62,115} \\
    \bottomrule
  \end{tabular}
\end{table}

\begin{table}[htbp]
  \centering
  \caption{储能设备在指定日期的实际充放电量（按四小时汇总，单位：kWh）}
  \label{tab:p3-storage}
  \footnotesize
  \begin{tabular}{lrrrrrr}
    \toprule
    时段 & \multicolumn{2}{c}{2025.3.20} & \multicolumn{2}{c}{2025.6.21} & \multicolumn{2}{c}{2025.9.23} \\
     & 充电 & 放电 & 充电 & 放电 & 充电 & 放电 \\
    \midrule
    0{:}00-4{:}00 & 4,183.11 & 436.34 & 604.87 & 1,022.39 & 4,716.66 & 324.79 \\
    4{:}00-8{:}00 & 907.49 & 5,900.81 & 469.00 & 4,396.07 & 1,018.72 & 6,581.69 \\
    8{:}00-12{:}00 & 4,695.33 & 2,406.36 & 9,512.78 & 0.00 & 5,294.08 & 1,812.49 \\
    12{:}00-16{:}00 & 6,285.80 & 254.72 & 0.00 & 0.00 & 5,892.83 & 465.86 \\
    16{:}00-20{:}00 & 9.23 & 5,667.02 & 168.09 & 5,393.94 & 99.06 & 5,870.41 \\
    20{:}00-24{:}00 & 5,562.75 & 3,026.13 & 6,097.43 & 2,504.35 & 5,325.43 & 2,811.44 \\
    \midrule
    $0{:}00$ 储电量 & \multicolumn{2}{c}{6,333.5} & \multicolumn{2}{c}{7,292.5} & \multicolumn{2}{c}{5,722.1} \\
    $24{:}00$ 储电量 & \multicolumn{2}{c}{6,155.7} & \multicolumn{2}{c}{7,663.1} & \multicolumn{2}{c}{5,982.4} \\
  \end{tabular}
\end{table}

\begin{table}[htbp]
  \centering
  \caption{2025 年 12 月 21 日储能设备的实际充放电量与 $0{:}00$、$24{:}00$ 储电量}
  \label{tab:p3-storage-dec}
  \footnotesize
  \begin{tabular}{lrr}
    \toprule
    时段 & 充电量/kWh & 放电量/kWh \\
    \midrule
    0{:}00-4{:}00 & 4,763.94 & 182.24 \\
    4{:}00-8{:}00 & 1,462.17 & 1,495.25 \\
    8{:}00-12{:}00 & 5,881.35 & 7,168.08 \\
    12{:}00-16{:}00 & 6,820.49 & 2,287.09 \\
    16{:}00-20{:}00 & 9.57 & 5,752.84 \\
    20{:}00-24{:}00 & 5,548.98 & 2,848.52 \\
    \midrule
    $0{:}00$ 储电量/kWh & \multicolumn{2}{c}{6,134.5} \\
    $24{:}00$ 储电量/kWh & \multicolumn{2}{c}{6,245.6} \\
    \bottomrule
  \end{tabular}
\end{table}

\begin{table}[htbp]
  \centering
  \caption{微网在指定日期的紧急购电事件（问题三，按同日连续时段合并）}
  \label{tab:p3-emg}
  \footnotesize
  \begin{tabular}{llrl}
    \toprule
    日期 & 紧急购电时间段 & 紧急购电量/kWh & 该日紧急费/元 \\
    \midrule
    \multirow{1}{*}{2025.3.20} & 9{:}00-10{:}00 & 371.02 & \multirow{1}{*}{1,819} \\
    \midrule
    2025.6.21 & \multicolumn{2}{c}{无} & 0.00 \\
    \multirow{2}{*}{2025.9.23} & 9{:}20-9{:}50 & 84.13 & \multirow{2}{*}{3,847} \\
     & 20{:}20-21{:}50 & 589.71 & \\
    \midrule
    \multirow{1}{*}{2025.12.21} & 7{:}50-8{:}00 & 12.14 & \multirow{1}{*}{71} \\
    \bottomrule
  \end{tabular}
  \par\vspace{2pt}
  \begin{minipage}{.92\textwidth}\footnotesize
    注：紧急购电量逐十分钟按 $5p_tr_t$ 计费后再把同日连续时段合并成事件，合并不改变费用；合并后电量与逐段之和严格相等。
  \end{minipage}
\end{table}

\begin{table}[htbp]
  \centering
  \caption{八种预报组合在 2025 年 2 月 1 日--12 月 31 日的实际费用（各组合单独标定参数，评价区间、负载模型、执行规则与计费口径完全一致）}
  \label{tab:p3-combo}
  \footnotesize
  \begin{tabular}{lrrrrrr}
    \toprule
    使用的预报时刻 & 初始计划费/元 & 调整费/元 & 紧急费/元 & 合计费用/元 & 相对 $\varnothing$ & 期末储电量/kWh \\
    \midrule
    仅 0:00 & 13,460,229 & 0 & 747,173 & \textbf{14,207,402} & --- & 6,290.4 \\
    {6} & 13,184,766 & 211,122 & 580,165 & \textbf{13,976,054} & -1.63\% & 6,290.4 \\
    {12} & 13,027,382 & 253,395 & 625,467 & \textbf{13,906,244} & -2.12\% & 6,057.1 \\
    {18} & 12,922,577 & 369,782 & 440,657 & \textbf{13,733,016} & -3.34\% & 6,057.1 \\
    {6,12} & 12,748,432 & 469,452 & 460,787 & \textbf{13,678,671} & -3.72\% & 5,891.2 \\
    {6,18} & 12,628,812 & 625,335 & 220,381 & \textbf{13,474,527} & -5.16\% & 5,891.2 \\
    {12,18} & 12,819,205 & 403,719 & 498,346 & \textbf{13,721,269} & -3.42\% & 6,057.1 \\
    {6,12,18} & 12,492,137 & 675,387 & 290,173 & \textbf{13,457,698} & -5.28\% & 5,891.2 \\
    \midrule
    对照：重新优化但仍用 0{:}00 版本 & --- & --- & --- & \textbf{13,658,321} & -3.86\% & 6,057.1 \\
    \midrule
    敏感性：退款口径（$x_t$ 可双向调整） & --- & --- & --- & \textbf{13,396,938} & --- & --- \\
    \bottomrule
  \end{tabular}
  \par\vspace{2pt}
  \begin{minipage}{.92\textwidth}\footnotesize
    注：所有方案都使用 0:00 发布的预报；组合列出的时刻是\textbf{额外}使用的版本。费用按主计费口径计算，\"合计费用\"已经包含初始计划费，不能与计划费、调整费再相加。期末储电量用于说明各方案的跨日状态差异，不能只比较费用。末行为计费口径的敏感性检查：退款口径的目标函数另含调减收益，其总额与上表的\"相对\"列不可直接比较。
  \end{minipage}
\end{table}


\subsection{求解结果与九项校验}

以 2025 年 1 月为预热期、$2$ 月 $1$ 日--$12$ 月 $31$ 日共 $334$ 天为正式区间，
主策略（使用全部四个发布时刻）的结果为 \textbf{合计购电费
$13{,}418{,}626.35$~元}。按性质拆开：计划购电费
$\sum_tp_tg_t^0=12{,}583{,}676.62$~元；上调增费
$\sum_t1.5p_tu_t=804{,}431.51$~元；调减退费
$\sum_t0.5p_tv_t=242{,}417.95$~元；两项相抵得\textbf{调整净额}
$562{,}013.56$~元，故\textbf{计划调整费合计} $13{,}145{,}690.17$~元；
\textbf{紧急购电费} $272{,}936.18$~元。

电量方面，初始计划量
$G^{\mathrm{plan}}=20{,}762{,}854.17$~kWh，最终生效常规量
$G^{\mathrm{eff}}=20{,}643{,}456.67$~kWh（其中上调 $656{,}263.05$~kWh、调减
$775{,}660.56$~kWh），紧急购电量 $G^{\mathrm{emg}}=47{,}350.77$~kWh，故微网
实际收到 $G^{\mathrm{recv}}=20{,}690{,}807.44$~kWh；储能充电
$6{,}689{,}851.97$~kWh、放电 $5{,}419{,}388.14$~kWh、弃电
$1{,}481{,}154.12$~kWh；评价区间首日 $0{:}00$ 储电量 $6{,}473.78$~kWh，末日期末
储电量 $5{,}798.18$~kWh。全年共正式提交调整 $334\times3=1002$ 次、求解线性规划
$334\times4=1336$ 次，与"四个交付块各求解一次"完全一致。

对主策略的 $334$ 个交付日逐日执行九项检查，结果见表~\ref{tab:p3-check}，
\textbf{九项的未通过天数全部为零}。数值上，逐段供需平衡的最大残差为
$2.27\times10^{-13}$~kWh，储能状态递推的最大绝对偏差与两项费用按独立路径复算的
偏差\textbf{恰为 $0$}；储电量上下界、充放电功率上限、同时充放电、购电量非负与跨日
状态衔接五项由模型约束与执行规则\textbf{结构性地}保证（充放电量在生成时就被逐段
截断到边界内），故逐日实测值恒为零。这些残差都在双精度浮点的舍入量级，说明模型
约束与账本记录是同一套代数关系的两种写法。

\subsection{结果分析}

\textbf{（1）总费用与紧急费沿预报格严格单调，计划调整费则不是。}把八种预报使用
方式的三项费用并列（表~\ref{tab:p3-combo}、图~\ref{fig:p3-combo}）可以看到：在
全部 $12$ 条"多用一个预报版本"的覆盖边上，\textbf{总费用严格不增}（由仅用
$0{:}00$ 的 $14{,}977{,}809$~元降到全部使用的 $13{,}418{,}626$~元，省
$1{,}559{,}183$~元、$10.41\%$），\textbf{紧急费也严格不增}（由 $2{,}382{,}673$
元降到 $272{,}936$~元，降幅 $88.5\%$）。多得到的预报并没有让总购电量凭空变少——
$G^{\mathrm{plan}}$ 在八种方式间只相差 $0.16\%$——而是把原本要在执行阶段按 $5$
倍电价紧急补购的电量，提前识别出来并按 $1.5$ 倍的价格调整；由于 $1.5p<5p$，这一
置换在总体上必然有利。\textbf{但计划调整费并不单调}：它在 $12$ 条覆盖边上有
$10$ 条上升，却在 $\{18\}\to\{12,18\}$（$13{,}187{,}014\to13{,}033{,}851$~元）
与 $\{6,18\}\to\{6,12,18\}$（$13{,}259{,}141\to13{,}145{,}690$~元）这两条边上
\textbf{反而下降}。原因是 $12{:}00$ 版重排了储能与购电，把一部分原本要靠上调补足
的缺口改由已经买好的电量覆盖，于是调整净额降低而总费用仍然下降。因此"版本越多越
省"是\textbf{总费用}层面的结论，不能简化为"调整费单项随版本单调增加"。

\begin{figure}[htbp]
  \centering
  \includegraphics[width=.72\textwidth]{p3_fig4_combos.png}
  \caption{八种预报使用方式的实际总费用分解（最左一项为只用 $0{:}00$ 预报并关闭
           全天调整的对照）}
  \label{fig:p3-combo}
\end{figure}

\textbf{（2）最值钱的版本并不是最准的那个。}为回答"是否需要其他时刻的预报"，
还要看\textbf{边际}改善 $\Delta_h(S)=C(S)-C(S\cup\{h\})$，即已经用了 $S$ 之后再
加用 $h$ 能省多少（完整表见附录第 \ref{app:p3-combo-detail} 节）。从零开始单独
引入时，$18{:}00$ 版省 $117.82$~万元、$12{:}00$ 版省 $94.59$~万元、$6{:}00$ 版
省 $49.78$~万元；并且 $18{:}00$ 版在 $\{6\}$、$\{12\}$、$\{6,12\}$ 之上仍然都是
三者中改善最大的（分别省 $94.69$、$45.17$、$46.23$~万元）。把这个排序与
表~\ref{tab:p3-skill} 的精度排序对照会发现\textbf{两者并不一致}：在白天窗口
$12{:}00$--$18{:}00$ 上，$12{:}00$ 版的平均绝对误差最小（$181.3$~kW），
$6{:}00$ 版 $304.9$~kW，$0{:}00$ 版 $445.8$~kW；而 $18{:}00$ 版\textbf{对白天没有
任何覆盖}，它只报 $18{:}00$--$24{:}00$ 这六小时。一个不覆盖白天、白天误差也无从
谈起的版本反而最值钱，说明它的价值\textbf{不在精度，而在它决定的是哪一个交付块}：
$18{:}00$ 版一旦提交，锁定的就是晚间 $18{:}00$--$24{:}00$——这既是电价高峰段，
又是日末回充的窗口，而且是全天唯一\textbf{没有光伏}的时段，风险全部落在没有外部
预报、只能靠同星期往期外推的\textbf{负载}上。反过来，$12{:}00$ 版的边际价值在
$\{6,18\}$ 之上只有 $11.44$~万元（$0.85\%$），远小于它单独引入时的 $94.59$~万元：
它对下午的修正 $6{:}00$ 版已经做过一次，$18{:}00$ 版又会覆盖相同的目标时段，信息
高度重复。

若必须精简，按"撤下的代价"排序：撤掉 $18{:}00$ 版多花 $46.23$~万元、撤掉
$6{:}00$ 版多花 $16.16$~万元、撤掉 $12{:}00$ 版只多花 $11.44$~万元。故本文的建议
是\textbf{三个时刻全用}；若要减少一次订阅与一次求解，可以省去 $12{:}00$ 版，代价
约为全年费用的 $0.85\%$。

\textbf{（3）口径与风险刻画的影响都远小于预报本身。}表~\ref{tab:p3-billing} 给出
两处口径的敏感性。计费口径之争（调减是否退款）在全年总费用上相差 $24.24$~万元
（$1.81\%$）：把主策略自身的物理量按不退款口径结算，计划调整费由
$13{,}145{,}690$~元升至 $13{,}388{,}108$~元；完全按不退款口径重新优化后为
$13{,}618{,}944$~元，比主口径高 $1.49\%$。两臂的物理运行差异很小，差额主要来自
\textbf{结算规则}而非策略。把"所有时段的历史误差合并成一个分位数"换成"每个时段
各自一个分位数"只差 $0.25\%$（$13{,}384{,}829$~元），说明风险刻画的\textbf{粒度}
不是关键；真正决定成本的是 $\alpha$ 的水平，见表~\ref{tab:p3-param}。

\textbf{（4）与问题二的对照。}问题二只能被动接受日前计划、缺口按 $5$ 倍电价硬补，
在同样的 $334$ 天上紧急购电量为 $140{,}301.24$~kWh、紧急费 $755{,}499.17$~元；
问题三主动付费调整，紧急购电量降到 $47{,}350.77$~kWh（$33.8\%$）、紧急费降到
$272{,}936.18$~元（$36.1\%$），合计购电费由 $13{,}939{,}720.46$~元降到
$13{,}418{,}626.35$~元，省 $521{,}094.11$~元（$3.74\%$）。需要说明的是，这一
节省\textbf{不是靠在年末掏空储能换来的}：问题二期末储电量为 $1{,}559.72$~kWh，
已经贴近下界；问题三日末储电量为 $5{,}798.18$~kWh，比前者高 $4{,}238$~kWh，
却仍然更省。日内三次发布的信息把大部分缺口从"$5p$ 的被动补购"搬到了"$1.5p$ 的
主动追加"，这正是第（1）条所测得的单调性的经济含义。

\textbf{（5）三个交付块的调整并不平均，储能轨迹会被逐次重新拼接。}$1002$ 次提交
中有相当一部分的结果是"保持不变"——全年被真正改动的时段占 $34.91\%$。四个指定
日期的逐时段结果见表~\ref{tab:p3-buy}、表~\ref{tab:p3-storage}、
表~\ref{tab:p3-storage-dec} 与表~\ref{tab:p3-emg}；调整在三个交付块间的分布、
储能轨迹的滚动修正与收益的月份差异，见附录第 \ref{app:p3-analysis} 节。

\subsection{模型的边界}

本文求解的是"风险余量 $+$ 滚动调整 $+$ 因果执行规则"的\textbf{可实现近似策略}，
每个决策时点只使用当时已经发布的信息，因而不是未知未来下的全局最优解；参数在
$31$ 天窗口内标定后冻结，是\textbf{近视最优}，且已如上文所示存在过拟合风险；
计费口径的歧义会改变调减的方向与账本。\textbf{五条适用边界的完整讨论见
附录~\ref{app:extra}}。
